\documentclass[11pt,a4paper]{article}

% ============================================================
% PACKAGES
% ============================================================
\usepackage[utf8]{inputenc}
\usepackage[T5]{fontenc}
\usepackage[vietnamese]{babel}

\usepackage{amsmath,amssymb,amsfonts}
\usepackage{bm}
\usepackage{booktabs}
\usepackage{multirow}
\usepackage{array}
\usepackage{graphicx}
\usepackage{xcolor}
\usepackage{enumitem}
\usepackage{algorithm}
\usepackage{algpseudocode}
\usepackage{tikz}
\usepackage{url}
\usepackage{hyperref}
\usepackage{float}
\usepackage{geometry}

\usetikzlibrary{
	arrows.meta,
	positioning,
	shapes.geometric,
	fit,
	calc
}

\geometry{
	left=2.5cm,
	right=2.5cm,
	top=2.5cm,
	bottom=2.5cm
}

% ============================================================
% DOCUMENT
% ============================================================

\title{
	\textbf{
		BeliefMove: Bayesian Belief-Augmented Large Language Model Agents
		for Uncertainty-Aware Next GPS Point Prediction
	}
}

\author{}
\date{}

\begin{document}
	
	\maketitle
	
	% ============================================================
	\begin{abstract}
		
		Dự đoán điểm GPS kế tiếp là một bài toán quan trọng trong khai phá
		trajectory, phân tích hành vi di chuyển và hệ thống giao thông thông minh.
		Các mô hình học sâu truyền thống có khả năng khai thác phụ thuộc
		không--thời gian nhưng thường gặp khó khăn khi trajectory thưa, nhiễu hoặc
		có nhiều hướng di chuyển khả dĩ. Gần đây, Large Language Models, LLMs,
		và các mô hình tác tử dựa trên LLM mở ra một hướng tiếp cận mới cho bài
		toán dự đoán vị trí thông qua khả năng suy luận ngữ nghĩa, khai thác tri
		thức và mô hình hóa ý định di chuyển. Tuy nhiên, LLM thường đưa ra một
		kết quả dự đoán duy nhất và chưa biểu diễn một cách có cấu trúc sự bất định
		giữa nhiều giả thuyết di chuyển.
		
		Bài nghiên cứu này đề xuất \textbf{BeliefMove}, một framework kết hợp
		LLM Agent với Bayesian Belief Network, BBN, cho dự đoán điểm GPS kế tiếp.
		Thay vì yêu cầu LLM trực tiếp dự đoán tọa độ kế tiếp, BeliefMove sử dụng
		LLM Agent để phân tích trajectory, suy luận ý định di chuyển và sinh ra
		một tập các giả thuyết về hướng hoặc tuyến đường tiếp theo. BBN sau đó
		chuyển các giả thuyết này thành một trạng thái belief xác suất và cập nhật
		trạng thái đó khi các quan sát GPS mới xuất hiện. Một cơ chế adaptive
		reasoning được đề xuất để cho phép Agent quyết định khi nào cần thực hiện
		thêm suy luận dựa trên entropy của belief state và mức độ bất đồng giữa
		LLM và BBN. Framework này hướng tới đồng thời cải thiện độ chính xác,
		khả năng chịu nhiễu, khả năng hiệu chỉnh xác suất và hiệu quả suy luận.
		Nghiên cứu được đánh giá trên nhiều tập dữ liệu GPS trajectory với các
		thí nghiệm về dự đoán điểm kế tiếp, dự đoán nhiều bước, dữ liệu thưa,
		nhiễu GPS, tổng quát hóa giữa các thành phố và calibration của độ tin cậy.
		
	\end{abstract}
	
	% ============================================================
	\section{Giới thiệu}
	
	GPS trajectory cung cấp một nguồn dữ liệu quan trọng để nghiên cứu hành
	vi di chuyển của con người và phương tiện. Một trajectory có thể được
	biểu diễn dưới dạng chuỗi các quan sát:
	
	\begin{equation}
		\mathcal{T}_{1:t}
		=
		\left\{
		p_1,p_2,\ldots,p_t
		\right\},
	\end{equation}
	
	trong đó mỗi điểm GPS được biểu diễn bởi:
	
	\begin{equation}
		p_i =
		(lat_i,lon_i,time_i).
	\end{equation}
	
	Một bài toán cơ bản trong trajectory mining là dự đoán điểm GPS tiếp theo:
	
	\begin{equation}
		p_{t+1}
		=
		(lat_{t+1},lon_{t+1}),
	\end{equation}
	
	dựa trên lịch sử trajectory:
	
	\begin{equation}
		P
		\left(
		p_{t+1}
		\mid
		p_1,\ldots,p_t
		\right).
	\end{equation}
	
	Các mô hình truyền thống như Markov model, Hidden Markov Model,
	Recurrent Neural Network, LSTM và Transformer đã được sử dụng rộng rãi
	cho bài toán này. Tuy nhiên, các mô hình này thường tập trung vào việc
	học quan hệ thống kê giữa các điểm GPS và chưa khai thác đầy đủ các yếu
	tố ngữ nghĩa như ý định di chuyển, mục tiêu chuyến đi, hoạt động của
	người dùng và tri thức về không gian.
	
	Sự phát triển của Large Language Models tạo ra một hướng tiếp cận mới.
	LLM có khả năng biểu diễn và suy luận trên thông tin không gian--thời
	gian ở dạng ngôn ngữ. Các nghiên cứu gần đây về LLM-based mobility agents,
	trong đó có AgentMove, cho thấy việc sử dụng memory, world knowledge và
	agentic reasoning có thể cải thiện khả năng dự đoán vị trí kế tiếp.
	
	Tuy nhiên, một hạn chế quan trọng vẫn tồn tại. Khi có nhiều hướng di
	chuyển khả dĩ, LLM thường phải lựa chọn một giả thuyết cuối cùng thay vì
	duy trì một phân phối xác suất trên các giả thuyết.
	
	Giả sử tại vị trí hiện tại có ba khả năng:
	
	\begin{equation}
		H =
		\left\{
		H_1,H_2,H_3
		\right\},
	\end{equation}
	
	trong đó:
	
	\begin{itemize}
		\item $H_1$: tiếp tục đi thẳng;
		\item $H_2$: rẽ phải;
		\item $H_3$: rẽ trái.
	\end{itemize}
	
	Một hệ thống dự đoán thông thường có thể lựa chọn:
	
	\begin{equation}
		\hat{H}
		=
		\arg\max_{H_i} P(H_i\mid\mathcal{T}_{1:t}).
	\end{equation}
	
	Tuy nhiên, việc chỉ giữ lại một giả thuyết sẽ làm mất thông tin về bất
	định.
	
	Nghiên cứu này đề xuất thay thế cách tiếp cận trên bằng một trạng thái
	belief:
	
	\begin{equation}
		B_t
		=
		\left\{
		P(H_1),
		P(H_2),
		\ldots,
		P(H_K)
		\right\}.
	\end{equation}
	
	Trạng thái belief được cập nhật liên tục khi quan sát mới xuất hiện:
	
	\begin{equation}
		B_t
		\rightarrow
		B_{t+1}.
	\end{equation}
	
	Ý tưởng trung tâm của nghiên cứu là kết hợp ba khả năng:
	
	\begin{equation}
		\boxed{
			\text{LLM Reasoning}
			+
			\text{Bayesian Belief}
			+
			\text{Agentic Decision Making}
		}
	\end{equation}
	
	Trong đó LLM chịu trách nhiệm suy luận ngữ nghĩa và sinh giả thuyết,
	Bayesian Belief Network định lượng và cập nhật bất định, còn Agent quyết
	định khi nào cần suy luận thêm.
	
	% ============================================================
	\section{Động lực nghiên cứu}
	
	\subsection{Hạn chế của dự đoán tất định}
	
	Một mô hình dự đoán thông thường có dạng:
	
	\begin{equation}
		\hat{p}_{t+1}
		=
		f
		\left(
		p_1,\ldots,p_t
		\right).
	\end{equation}
	
	Cách biểu diễn này ngầm giả định rằng có một kết quả tối ưu duy nhất.
	Tuy nhiên, mobility trajectory thường có tính đa phương thức.
	
	Từ cùng một vị trí, phương tiện có thể:
	
	\begin{equation}
		P(p_{t+1}\mid\mathcal{T})
		=
		\sum_{k=1}^{K}
		P(p_{t+1}\mid H_k,\mathcal{T})
		P(H_k\mid\mathcal{T}).
	\end{equation}
	
	Do đó, việc mô hình hóa phân phối trên các giả thuyết là phù hợp hơn
	với bản chất của bài toán.
	
	\subsection{LLM có khả năng reasoning nhưng chưa có belief state rõ ràng}
	
	LLM có thể suy luận rằng một trajectory có khả năng tương ứng với hành
	trình đi làm:
	
	\begin{equation}
		P(
		Intent=\text{commuting}
		\mid
		\mathcal{T}
		)
		\approx 0.72.
	\end{equation}
	
	Tuy nhiên, confidence do LLM sinh ra không nhất thiết là một xác suất
	được hiệu chỉnh tốt.
	
	Vì vậy, nghiên cứu đề xuất tách hai vai trò:
	
	\begin{equation}
		\text{LLM}
		\rightarrow
		\text{Semantic Hypothesis}
	\end{equation}
	
	và:
	
	\begin{equation}
		\text{BBN}
		\rightarrow
		\text{Probabilistic Belief}.
	\end{equation}
	
	\subsection{Động lực cho adaptive reasoning}
	
	Không phải mọi trajectory đều cần nhiều bước suy luận.
	
	Nếu:
	
	\begin{equation}
		H(B_t)
		=
		-\sum_i P(H_i)\log P(H_i)
	\end{equation}
	
	nhỏ, Agent có thể trực tiếp dự đoán.
	
	Ngược lại, nếu:
	
	\begin{equation}
		H(B_t)>\tau,
	\end{equation}
	
	Agent cần thực hiện thêm reasoning hoặc khai thác thêm context.
	
	Điều này cho phép giảm chi phí gọi LLM trong các trường hợp prediction
	đã có độ tin cậy cao.
	
	% ============================================================
	\section{Khoảng trống nghiên cứu}
	
	Nghiên cứu đề xuất tập trung vào bốn khoảng trống chính.
	
	\begin{enumerate}
		
		\item
		Các mô hình LLM-based mobility reasoning có khả năng suy luận ngữ nghĩa
		nhưng chưa khai thác đầy đủ một belief state xác suất có thể cập nhật
		theo thời gian.
		
		\item
		Các mô hình Bayesian truyền thống có khả năng biểu diễn bất định nhưng
		khó khai thác world knowledge và semantic mobility intent ở mức cao.
		
		\item
		Hai loại reasoning trên thường được nghiên cứu riêng biệt thay vì kết
		hợp trong một agentic framework.
		
		\item
		Chi phí suy luận của LLM có thể cao nếu Agent luôn thực hiện nhiều bước
		reasoning, trong khi không phải mọi trạng thái trajectory đều có mức độ
		bất định cao.
		
	\end{enumerate}
	
	Từ đó, nghiên cứu đặt ra câu hỏi:
	
	\begin{quote}
		Liệu việc kết hợp semantic reasoning của LLM với Bayesian belief reasoning
		và adaptive agentic decision making có thể cải thiện độ chính xác,
		độ tin cậy và khả năng chống chịu của dự đoán điểm GPS kế tiếp hay không?
	\end{quote}
	
	% ============================================================
	\section{Mục tiêu nghiên cứu}
	
	Mục tiêu tổng quát là xây dựng một framework LLM Agent có khả năng dự đoán
	điểm GPS kế tiếp trong điều kiện bất định.
	
	Các mục tiêu cụ thể gồm:
	
	\begin{enumerate}
		
		\item
		Xây dựng biểu diễn trajectory phù hợp cho LLM-based mobility reasoning.
		
		\item
		Thiết kế Mobility Belief Network biểu diễn các biến ẩn như mobility
		intent, destination và next road.
		
		\item
		Thiết kế cơ chế để LLM Agent sinh các mobility hypotheses.
		
		\item
		Thiết kế Bayesian belief updating khi quan sát GPS mới.
		
		\item
		Thiết kế cơ chế adaptive reasoning dựa trên uncertainty.
		
		\item
		Đánh giá khả năng dự đoán trong điều kiện GPS thưa và nhiễu.
		
		\item
		Đánh giá uncertainty calibration.
		
		\item
		Đánh giá khả năng giảm chi phí suy luận LLM.
		
	\end{enumerate}
	
	% ============================================================
	\section{Câu hỏi nghiên cứu}
	
	\subsection{RQ1}
	
	Liệu Bayesian belief reasoning có cải thiện độ chính xác dự đoán điểm
	GPS kế tiếp so với LLM Agent thuần túy hay không?
	
	\subsection{RQ2}
	
	Liệu belief updating có giúp mô hình chống chịu tốt hơn với GPS bị thiếu
	hoặc nhiễu hay không?
	
	\subsection{RQ3}
	
	Liệu Bayesian belief có giúp confidence của mô hình được calibration
	tốt hơn hay không?
	
	\subsection{RQ4}
	
	Liệu adaptive reasoning có thể giảm số lần gọi LLM mà không làm giảm
	đáng kể độ chính xác hay không?
	
	% ============================================================
	\section{Framework BeliefMove}
	
	\subsection{Tổng quan}
	
	Framework đề xuất có sáu thành phần:
	
	\begin{enumerate}
		\item GPS trajectory parser;
		\item Spatial-temporal trajectory encoder;
		\item LLM Mobility Agent;
		\item Mobility Hypothesis Generator;
		\item Bayesian Belief Network;
		\item Adaptive Reasoning Controller.
	\end{enumerate}
	
	Kiến trúc tổng quát được mô tả trong Hình~\ref{fig:architecture}.
	
	\begin{figure}[H]
		\centering
		
		\begin{tikzpicture}[
			node distance=8mm and 10mm,
			block/.style={
				rectangle,
				rounded corners,
				draw,
				align=center,
				minimum width=32mm,
				minimum height=10mm
			},
			decision/.style={
				diamond,
				draw,
				align=center,
				aspect=2,
				inner sep=2pt
			},
			arrow/.style={
				-{Latex[length=2mm]}
			}
			]
			
			\node[block] (gps) {GPS\\Trajectory};
			
			\node[block, below=of gps] (encoder)
			{Spatial--Temporal\\Encoder};
			
			\node[block, below=of encoder] (agent)
			{LLM Mobility\\Agent};
			
			\node[block, below left=12mm and 8mm of agent] (hyp)
			{Mobility\\Hypotheses};
			
			\node[block, below right=12mm and 8mm of agent] (context)
			{World and\\Context Knowledge};
			
			\node[block, below=25mm of agent] (bbn)
			{Bayesian Belief\\Network};
			
			\node[decision, below=of bbn] (decision)
			{Uncertainty\\low?};
			
			\node[block, below left=10mm and 8mm of decision] (predict)
			{Next GPS\\Prediction};
			
			\node[block, below right=10mm and 8mm of decision] (reason)
			{Additional\\Reasoning};
			
			\node[block, below=of reason] (update)
			{Belief\\Update};
			
			\draw[arrow] (gps) -- (encoder);
			\draw[arrow] (encoder) -- (agent);
			
			\draw[arrow] (agent) -- (hyp);
			\draw[arrow] (agent) -- (context);
			
			\draw[arrow] (hyp) -- (bbn);
			\draw[arrow] (context) -- (bbn);
			
			\draw[arrow] (bbn) -- (decision);
			
			\draw[arrow] (decision) -- node[above left]{Yes} (predict);
			\draw[arrow] (decision) -- node[above right]{No} (reason);
			
			\draw[arrow] (reason) -- (update);
			\draw[arrow] (update) |- (bbn);
			
		\end{tikzpicture}
		
		\caption{
			Kiến trúc tổng thể của BeliefMove.
		}
		\label{fig:architecture}
		
	\end{figure}
	
	% ============================================================
	\section{Biểu diễn bài toán}
	
	Cho trajectory:
	
	\begin{equation}
		\mathcal{T}_{1:t}
		=
		\{p_1,p_2,\ldots,p_t\}.
	\end{equation}
	
	Mỗi điểm:
	
	\begin{equation}
		p_i
		=
		(
		lat_i,
		lon_i,
		time_i,
		v_i,
		h_i
		),
	\end{equation}
	
	trong đó:
	
	\begin{itemize}
		\item $lat_i$: latitude;
		\item $lon_i$: longitude;
		\item $time_i$: timestamp;
		\item $v_i$: tốc độ;
		\item $h_i$: hướng di chuyển.
	\end{itemize}
	
	Ngoài trajectory, hệ thống có thể sử dụng context:
	
	\begin{equation}
		C_t
		=
		\{
		road_t,
		region_t,
		POI_t,
		traffic_t,
		weather_t
		\}.
	\end{equation}
	
	Mục tiêu là học:
	
	\begin{equation}
		P
		\left(
		p_{t+1}
		\mid
		\mathcal{T}_{1:t},
		C_t
		\right).
	\end{equation}
	
	% ============================================================
	\section{Mobility State}
	
	Thay vì đưa trực tiếp latitude và longitude vào BBN, nghiên cứu định
	nghĩa một mobility state:
	
	\begin{equation}
		S_t
		=
		\{
		R_t,
		A_t,
		I_t,
		D_t,
		V_t,
		H_t
		\},
	\end{equation}
	
	trong đó:
	
	\begin{itemize}
		\item $R_t$: road segment;
		\item $A_t$: activity;
		\item $I_t$: mobility intent;
		\item $D_t$: destination;
		\item $V_t$: speed state;
		\item $H_t$: heading state.
	\end{itemize}
	
	Điểm GPS kế tiếp được suy ra từ mobility state:
	
	\begin{equation}
		P(p_{t+1}\mid\mathcal{T}_{1:t})
		=
		\sum_{S_t}
		P(p_{t+1}\mid S_t)
		P(S_t\mid\mathcal{T}_{1:t}).
	\end{equation}
	
	% ============================================================
	\section{LLM Mobility Agent}
	
	LLM Agent nhận trajectory đã được chuẩn hóa thành một biểu diễn có ý
	nghĩa:
	
	\begin{verbatim}
		time = 08:10
		road = R125
		speed = 31 km/h
		heading = East
		region = District_1
		
		time = 08:11
		road = R126
		speed = 34 km/h
		heading = East
		region = District_1
	\end{verbatim}
	
	Agent thực hiện ba nhiệm vụ.
	
	\subsection{Trajectory interpretation}
	
	Agent phân tích:
	
	\begin{equation}
		Z_t
		=
		f_{\mathrm{LLM}}
		(
		\mathcal{T}_{1:t},
		C_t
		).
	\end{equation}
	
	\subsection{Intent inference}
	
	Agent suy luận:
	
	\begin{equation}
		P_{LLM}(I_t\mid Z_t).
	\end{equation}
	
	Ví dụ:
	
	\begin{verbatim}
		commuting = 0.72
		shopping  = 0.18
		other     = 0.10
	\end{verbatim}
	
	\subsection{Hypothesis generation}
	
	Agent sinh:
	
	\begin{equation}
		\mathcal{H}_t
		=
		\{
		H_1,H_2,\ldots,H_K
		\}.
	\end{equation}
	
	Mỗi hypothesis có thể chứa:
	
	\begin{equation}
		H_k
		=
		(
		intent,
		destination,
		next\_road,
		next\_region
		).
	\end{equation}
	
	% ============================================================
	\section{Mobility Bayesian Belief Network}
	
	\subsection{Mục tiêu}
	
	BBN được sử dụng để mô hình hóa quan hệ xác suất giữa các biến mobility.
	
	Một mạng đề xuất:
	
	\begin{figure}[H]
		\centering
		
		\begin{tikzpicture}[
			node distance=12mm and 13mm,
			state/.style={
				rectangle,
				rounded corners,
				draw,
				minimum width=25mm,
				minimum height=8mm,
				align=center
			},
			arrow/.style={
				-{Latex[length=2mm]}
			}
			]
			
			\node[state] (time) {Time};
			
			\node[state, below=of time] (activity) {Activity};
			
			\node[state, left=of activity] (history) {Trajectory\\History};
			
			\node[state, right=of activity] (poi) {POI /\\Context};
			
			\node[state, below=of activity] (intent) {Mobility\\Intent};
			
			\node[state, below=of intent] (destination)
			{Destination};
			
			\node[state, left=of destination] (current)
			{Current\\Road};
			
			\node[state, right=of destination] (heading)
			{Heading};
			
			\node[state, below=of destination] (nextroad)
			{Next\\Road};
			
			\node[state, below=of nextroad] (region)
			{Next\\Region};
			
			\node[state, below=of region] (gps)
			{Next GPS\\Point};
			
			\draw[arrow] (time) -- (activity);
			\draw[arrow] (history) -- (intent);
			\draw[arrow] (activity) -- (intent);
			\draw[arrow] (poi) -- (intent);
			
			\draw[arrow] (intent) -- (destination);
			
			\draw[arrow] (current) -- (nextroad);
			\draw[arrow] (heading) -- (nextroad);
			\draw[arrow] (destination) -- (nextroad);
			
			\draw[arrow] (nextroad) -- (region);
			\draw[arrow] (region) -- (gps);
			
		\end{tikzpicture}
		
		\caption{
			Mobility Bayesian Belief Network.
		}
		\label{fig:bbn}
	\end{figure}
	
	Các quan hệ xác suất chính gồm:
	
	\begin{equation}
		P(I_t\mid History,Time,Activity,POI),
	\end{equation}
	
	\begin{equation}
		P(D_t\mid I_t,History),
	\end{equation}
	
	\begin{equation}
		P(R_{t+1}\mid R_t,H_t,D_t),
	\end{equation}
	
	và:
	
	\begin{equation}
		P(p_{t+1}\mid R_{t+1},V_t,H_t).
	\end{equation}
	
	% ============================================================
	\section{Bayesian Belief Updating}
	
	Tại thời điểm $t$, Agent duy trì belief:
	
	\begin{equation}
		B_t(H_i)
		=
		P(H_i\mid E_{1:t}),
	\end{equation}
	
	trong đó $E_{1:t}$ là toàn bộ evidence đã quan sát.
	
	Khi GPS mới $E_{t+1}$ xuất hiện:
	
	\begin{equation}
		B_{t+1}(H_i)
		=
		P(H_i\mid E_{1:t+1}).
	\end{equation}
	
	Theo Bayes:
	
	\begin{equation}
		B_{t+1}(H_i)
		=
		\frac{
			P(E_{t+1}\mid H_i)
			B_t(H_i)
		}{
			\sum_{j=1}^{K}
			P(E_{t+1}\mid H_j)
			B_t(H_j)
		}.
	\end{equation}
	
	Ví dụ ban đầu:
	
	\begin{equation}
		B_t
		=
		\{
		0.50,
		0.30,
		0.20
		\}.
	\end{equation}
	
	Sau quan sát mới:
	
	\begin{equation}
		B_{t+1}
		=
		\{
		0.68,
		0.22,
		0.10
		\}.
	\end{equation}
	
	Belief state vì vậy được cập nhật thay vì tái khởi tạo sau mỗi quan sát.
	
	% ============================================================
	\section{LLM--BBN Belief Fusion}
	
	LLM tạo:
	
	\begin{equation}
		P_{LLM}(H_i).
	\end{equation}
	
	BBN tạo:
	
	\begin{equation}
		P_{BBN}(H_i).
	\end{equation}
	
	Một baseline fusion đơn giản là:
	
	\begin{equation}
		P^*(H_i)
		=
		\alpha
		P_{LLM}(H_i)
		+
		(1-\alpha)
		P_{BBN}(H_i).
	\end{equation}
	
	Tuy nhiên, nghiên cứu đề xuất sử dụng trọng số thích nghi:
	
	\begin{equation}
		\alpha_t
		=
		f
		\left(
		U_{LLM},
		U_{BBN},
		D_{LLM,BBN}
		\right),
	\end{equation}
	
	trong đó:
	
	\begin{equation}
		U_{LLM}
	\end{equation}
	
	là uncertainty của LLM,
	
	\begin{equation}
		U_{BBN}
	\end{equation}
	
	là uncertainty của BBN, và:
	
	\begin{equation}
		D_{LLM,BBN}
		=
		D_{KL}
		\left(
		P_{LLM}
		\|
		P_{BBN}
		\right)
	\end{equation}
	
	đo mức độ bất đồng giữa hai nguồn reasoning.
	
	Belief cuối cùng:
	
	\begin{equation}
		P^*(H_i)
		=
		\alpha_tP_{LLM}(H_i)
		+
		(1-\alpha_t)P_{BBN}(H_i).
	\end{equation}
	
	% ============================================================
	\section{Belief Memory}
	
	Một thành phần mới được đề xuất là \textbf{Bayesian Belief Memory}.
	
	Thay vì chỉ lưu text hoặc historical trajectory, Agent lưu:
	
	\begin{equation}
		M_t^{B}
		=
		\{
		B_1,B_2,\ldots,B_t
		\}.
	\end{equation}
	
	Mỗi belief state:
	
	\begin{equation}
		B_t
		=
		\{
		P(I_t),
		P(D_t),
		P(R_{t+1}),
		P(Region_{t+1})
		\}.
	\end{equation}
	
	Belief Memory cho phép Agent truy hồi các trạng thái bất định trước đó
	và cập nhật chúng khi trajectory mới xuất hiện.
	
	% ============================================================
	\section{Adaptive Uncertainty-Aware Reasoning}
	
	Entropy của belief được định nghĩa:
	
	\begin{equation}
		\mathcal{H}(B_t)
		=
		-\sum_{i=1}^{K}
		P(H_i)
		\log P(H_i).
	\end{equation}
	
	Nếu:
	
	\begin{equation}
		\mathcal{H}(B_t)<\tau,
	\end{equation}
	
	Agent thực hiện prediction.
	
	Nếu:
	
	\begin{equation}
		\mathcal{H}(B_t)\geq\tau,
	\end{equation}
	
	Agent thực hiện additional reasoning.
	
	Một policy đơn giản:
	
	\begin{equation}
		a_t
		=
		\begin{cases}
			predict,
			&
			\mathcal{H}(B_t)<\tau
			\\
			reason,
			&
			\mathcal{H}(B_t)\geq\tau.
		\end{cases}
	\end{equation}
	
	Một policy mở rộng có thể sử dụng:
	
	\begin{equation}
		a_t
		=
		\pi
		(
		B_t,
		\mathcal{H}(B_t),
		D_{LLM,BBN},
		C_t
		).
	\end{equation}
	
	Các action có thể gồm:
	
	\begin{equation}
		a_t
		\in
		\{
		predict,
		retrieve,
		reason,
		update,
		observe
		\}.
	\end{equation}
	
	% ============================================================
	\section{Dự đoán điểm GPS kế tiếp}
	
	Sau khi có belief:
	
	\begin{equation}
		P(R_{t+1}\mid E_{1:t}),
	\end{equation}
	
	mô hình dự đoán GPS:
	
	\begin{equation}
		P(
		p_{t+1}
		\mid
		R_{t+1},
		V_t,
		H_t,
		B_t
		).
	\end{equation}
	
	Điểm dự đoán có thể được tính:
	
	\begin{equation}
		\hat{p}_{t+1}
		=
		\sum_{r}
		P(R_{t+1}=r\mid B_t)
		\hat{p}_{t+1}^{(r)}.
	\end{equation}
	
	Trong trường hợp multi-modal prediction:
	
	\begin{equation}
		\mathcal{P}_{t+1}
		=
		\{
		(
		\hat{p}_{t+1}^{(1)},q_1),
		\ldots,
		(
		\hat{p}_{t+1}^{(K)},q_K)
		\}.
	\end{equation}
	
	Trong đó $q_i$ là xác suất của hypothesis tương ứng.
	
	% ============================================================
	\section{Multi-step Trajectory Forecasting}
	
	Framework có thể được mở rộng từ next-point prediction sang multi-step
	forecasting:
	
	\begin{equation}
		\hat{\mathcal{T}}_{t+1:t+K}
		=
		\{
		\hat{p}_{t+1},
		\ldots,
		\hat{p}_{t+K}
		\}.
	\end{equation}
	
	Sau mỗi bước, belief được cập nhật:
	
	\begin{equation}
		B_t
		\rightarrow
		B_{t+1}
		\rightarrow
		\cdots
		\rightarrow
		B_{t+K}.
	\end{equation}
	
	Do đó framework không chỉ dự đoán một điểm độc lập mà thực hiện
	sequential probabilistic forecasting.
	
	% ============================================================
	\section{Thuật toán BeliefMove}
	
	\begin{algorithm}[H]
		\caption{BeliefMove: Bayesian Belief-Augmented Mobility Agent}
		\label{alg:beliefmove}
		
		\begin{algorithmic}[1]
			
			\Require
			Trajectory $\mathcal{T}_{1:t}$,
			context $C_t$,
			belief state $B_t$,
			uncertainty threshold $\tau$
			
			\Ensure
			Next GPS prediction $\hat{p}_{t+1}$
			
			\State
			$Z_t
			\gets
			\mathrm{Encode}
			(\mathcal{T}_{1:t},C_t)$
			
			\State
			$S_t
			\gets
			\mathrm{LLMReason}
			(Z_t)$
			
			\State
			$\mathcal{H}_t
			\gets
			\mathrm{GenerateHypotheses}
			(S_t)$
			
			\State
			$P_{LLM}
			\gets
			\mathrm{EstimateLLMBelief}
			(\mathcal{H}_t)$
			
			\State
			$P_{BBN}
			\gets
			\mathrm{BayesianInference}
			(B_t,S_t)$
			
			\State
			$D
			\gets
			D_{KL}(P_{LLM}\|P_{BBN})$
			
			\State
			$B_t
			\gets
			\mathrm{AdaptiveFusion}
			(P_{LLM},P_{BBN},D)$
			
			\State
			$U
			\gets
			\mathrm{Entropy}(B_t)$
			
			\If{$U \geq \tau$}
			
			\State
			$C_t'
			\gets
			\mathrm{RetrieveAdditionalContext}(C_t)$
			
			\State
			$S_t'
			\gets
			\mathrm{AdditionalReasoning}(Z_t,C_t')$
			
			\State
			$B_t
			\gets
			\mathrm{UpdateBelief}(B_t,S_t')$
			
			\EndIf
			
			\State
			$\hat{p}_{t+1}
			\gets
			\mathrm{GPSDecoder}(B_t)$
			
			\State
			$B_{t+1}
			\gets
			\mathrm{BeliefUpdate}(B_t,p_{t+1})$
			
			\State
			\Return
			$\hat{p}_{t+1}$
			
		\end{algorithmic}
	\end{algorithm}
	
	% ============================================================
	\section{Thiết kế thực nghiệm}
	
	\subsection{Dataset}
	
	Nghiên cứu đề xuất sử dụng ít nhất ba tập dữ liệu:
	
	\begin{itemize}
		\item GeoLife;
		\item T-Drive;
		\item Porto Taxi.
	\end{itemize}
	
	Các dataset này đại diện cho các dạng mobility khác nhau, bao gồm
	human mobility và vehicle mobility.
	
	% ============================================================
	\subsection{Tiền xử lý}
	
	Pipeline:
	
	\begin{equation}
		Raw\ GPS
		\rightarrow
		Cleaning
		\rightarrow
		Segmentation
		\rightarrow
		Map\ Matching
		\rightarrow
		Spatial\ Representation
		\rightarrow
		Trajectory\ Sequence.
	\end{equation}
	
	Các bước gồm:
	
	\begin{enumerate}
		
		\item
		Sắp xếp dữ liệu theo user hoặc vehicle và timestamp.
		
		\item
		Loại bỏ các điểm GPS không hợp lệ.
		
		\item
		Tính tốc độ:
		
		\begin{equation}
			v_i
			=
			\frac{
				d(p_i,p_{i-1})
			}{
				t_i-t_{i-1}
			}.
		\end{equation}
		
		\item
		Loại bỏ hoặc đánh dấu các điểm có tốc độ bất thường.
		
		\item
		Map matching sang road segment nếu bản đồ đường được sử dụng.
		
		\item
		Chuyển trajectory thành các cửa sổ quan sát.
		
	\end{enumerate}
	
	% ============================================================
	\section{Các baseline}
	
	Các baseline được chia thành bốn nhóm.
	
	\subsection{Traditional models}
	
	\begin{enumerate}
		\item Historical Average;
		\item Markov Model;
		\item Kalman Filter.
	\end{enumerate}
	
	\subsection{Deep learning models}
	
	\begin{enumerate}
		\item LSTM;
		\item GRU;
		\item Transformer.
	\end{enumerate}
	
	\subsection{LLM-based models}
	
	\begin{enumerate}
		\item Direct LLM;
		\item AgentMove;
		\item Các LLM-based trajectory models phù hợp.
	\end{enumerate}
	
	\subsection{Proposed models}
	
	\begin{enumerate}
		\item AgentMove + BBN;
		\item AgentMove + Intent + BBN;
		\item BeliefMove;
		\item BeliefMove + Adaptive Reasoning.
	\end{enumerate}
	
	% ============================================================
	\section{Ablation Study}
	
	Ablation study gồm:
	
	\begin{table}[H]
		\centering
		\caption{Các phiên bản ablation của framework.}
		\label{tab:ablation}
		
		\begin{tabular}{lccccc}
			\toprule
			Model & LLM & Agent & Intent & BBN & Adaptive \\
			\midrule
			Transformer & & & & & \\
			LLM & \checkmark & & & & \\
			AgentMove & \checkmark & \checkmark & \checkmark & & \\
			AgentMove-BBN & \checkmark & \checkmark & & \checkmark & \\
			AgentMove-Intent-BBN & \checkmark & \checkmark & \checkmark & \checkmark & \\
			BeliefMove & \checkmark & \checkmark & \checkmark & \checkmark & \checkmark \\
			\bottomrule
		\end{tabular}
		
	\end{table}
	
	Mục tiêu của ablation là xác định đóng góp độc lập của từng thành phần.
	
	% ============================================================
	\section{Metrics}
	
	\subsection{Average Displacement Error}
	
	\begin{equation}
		ADE
		=
		\frac{1}{N}
		\sum_{i=1}^{N}
		d(
		p_i,\hat{p}_i
		).
	\end{equation}
	
	\subsection{Final Displacement Error}
	
	\begin{equation}
		FDE
		=
		d(
		p_T,\hat{p}_T
		).
	\end{equation}
	
	\subsection{Haversine Distance}
	
	Khoảng cách địa lý được tính bằng công thức Haversine.
	
	\begin{equation}
		d
		=
		2R
		\arcsin
		\left(
		\sqrt{
			\sin^2
			\frac{\Delta\phi}{2}
			+
			\cos(\phi_1)
			\cos(\phi_2)
			\sin^2
			\frac{\Delta\lambda}{2}
		}
		\right).
	\end{equation}
	
	\subsection{Top-$K$ Accuracy}
	
	\begin{equation}
		Accuracy@K
		=
		\frac{
			\text{Number of correct predictions within top-}K
		}{
			N
		}.
	\end{equation}
	
	% ============================================================
	\section{Uncertainty Calibration}
	
	Do framework sinh ra probability distribution, cần đánh giá calibration.
	
	\subsection{Negative Log-Likelihood}
	
	\begin{equation}
		NLL
		=
		-\frac{1}{N}
		\sum_{i=1}^{N}
		\log
		P(y_i\mid x_i).
	\end{equation}
	
	\subsection{Brier Score}
	
	\begin{equation}
		BS
		=
		\frac{1}{N}
		\sum_{i=1}^{N}
		(p_i-y_i)^2.
	\end{equation}
	
	\subsection{Expected Calibration Error}
	
	\begin{equation}
		ECE
		=
		\sum_{m=1}^{M}
		\frac{|B_m|}{N}
		\left|
		acc(B_m)-conf(B_m)
		\right|.
	\end{equation}
	
	% ============================================================
	\section{Thí nghiệm với GPS bị thiếu}
	
	Để đánh giá robustness, ngẫu nhiên loại bỏ một tỷ lệ GPS points:
	
	\begin{equation}
		r
		\in
		\{
		0,
		0.1,
		0.2,
		0.3,
		0.5
		\}.
	\end{equation}
	
	So sánh:
	
	\begin{equation}
		ADE(r),
		FDE(r).
	\end{equation}
	
	Giả thuyết:
	
	\begin{equation}
		\frac{\partial ADE}{\partial r}
		\Bigg|_{\text{BeliefMove}}
		<
		\frac{\partial ADE}{\partial r}
		\Bigg|_{\text{AgentMove}}.
	\end{equation}
	
	Điều này cho thấy belief state giúp mô hình duy trì thông tin khi quan sát
	bị thiếu.
	
	% ============================================================
	\section{Thí nghiệm với GPS nhiễu}
	
	GPS noise được mô hình hóa:
	
	\begin{equation}
		p_i'
		=
		p_i+\epsilon_i,
	\end{equation}
	
	trong đó:
	
	\begin{equation}
		\epsilon_i
		\sim
		\mathcal{N}(0,\sigma^2).
	\end{equation}
	
	Thử nghiệm với:
	
	\begin{equation}
		\sigma
		\in
		\{
		5m,
		10m,
		20m,
		50m,
		100m
		\}.
	\end{equation}
	
	% ============================================================
	\section{Cross-City Generalization}
	
	Mô hình được train trên một thành phố và test trên thành phố khác.
	
	Ví dụ:
	
	\begin{equation}
		Train=Beijing,
		\qquad
		Test=Porto.
	\end{equation}
	
	Đánh giá:
	
	\begin{equation}
		\Delta ADE
		=
		ADE_{cross-city}
		-
		ADE_{in-city}.
	\end{equation}
	
	Một mô hình tổng quát tốt nên có:
	
	\begin{equation}
		\Delta ADE
		\rightarrow
		0.
	\end{equation}
	
	% ============================================================
	\section{LLM Cost Analysis}
	
	Đánh giá:
	
	\begin{equation}
		Cost
		=
		f(
		N_{calls},
		N_{tokens},
		Latency
		).
	\end{equation}
	
	So sánh:
	
	\begin{table}[H]
		\centering
		\caption{Đánh giá hiệu quả suy luận.}
		\label{tab:efficiency}
		
		\begin{tabular}{lcccc}
			\toprule
			Model & LLM Calls & Tokens & Latency & ADE \\
			\midrule
			LLM & & & & \\
			AgentMove & & & & \\
			AgentMove-BBN & & & & \\
			BeliefMove & & & & \\
			\bottomrule
		\end{tabular}
		
	\end{table}
	
	Mục tiêu là chứng minh:
	
	\begin{equation}
		Cost_{BeliefMove}
		<
		Cost_{AgentMove}
	\end{equation}
	
	trong khi:
	
	\begin{equation}
		ADE_{BeliefMove}
		\leq
		ADE_{AgentMove}.
	\end{equation}
	
	% ============================================================
	\section{Nghiên cứu sự bất đồng giữa LLM và BBN}
	
	Định nghĩa:
	
	\begin{equation}
		D_t
		=
		D_{KL}
		\left(
		P_{LLM}
		\|
		P_{BBN}
		\right).
	\end{equation}
	
	Nếu:
	
	\begin{equation}
		D_t>\delta,
	\end{equation}
	
	Agent thực hiện reflection.
	
	Quy trình:
	
	\begin{equation}
		LLM
		\rightarrow
		BBN
		\rightarrow
		Disagreement
		\rightarrow
		Reflection
		\rightarrow
		Belief\ Update.
	\end{equation}
	
	Đây là một thí nghiệm quan trọng để chứng minh Agent không chỉ là một
	pipeline tuần tự mà có khả năng tự điều chỉnh reasoning.
	
	% ============================================================
	\section{Các giả thuyết nghiên cứu}
	
	\subsection{H1}
	
	BeliefMove có ADE và FDE thấp hơn AgentMove.
	
	\begin{equation}
		ADE_{BeliefMove}
		<
		ADE_{AgentMove}.
	\end{equation}
	
	\subsection{H2}
	
	BeliefMove có robustness tốt hơn khi GPS bị thiếu.
	
	\begin{equation}
		ADE_{BeliefMove}(r)
		<
		ADE_{AgentMove}(r).
	\end{equation}
	
	\subsection{H3}
	
	BeliefMove có calibration tốt hơn.
	
	\begin{equation}
		ECE_{BeliefMove}
		<
		ECE_{AgentMove}.
	\end{equation}
	
	\subsection{H4}
	
	Adaptive reasoning giảm chi phí LLM.
	
	\begin{equation}
		Calls_{Adaptive}
		<
		Calls_{FullReasoning}.
	\end{equation}
	
	% ============================================================
	\section{Các đóng góp dự kiến}
	
	Nghiên cứu dự kiến có bốn đóng góp chính.
	
	\begin{enumerate}
		
		\item
		Đề xuất BeliefMove, một framework kết hợp LLM Agent và Bayesian Belief
		Network cho dự đoán điểm GPS kế tiếp theo hướng probabilistic mobility
		reasoning.
		
		\item
		Đề xuất Mobility Belief Network biểu diễn có cấu trúc các biến ẩn gồm
		mobility intent, destination, next road và next region.
		
		\item
		Đề xuất Bayesian Belief Memory và adaptive belief updating để duy trì
		và cập nhật nhiều mobility hypotheses trong quá trình trajectory tiếp
		diễn.
		
		\item
		Đề xuất adaptive uncertainty-aware reasoning dựa trên entropy và
		LLM--BBN disagreement nhằm cân bằng prediction accuracy và LLM inference
		cost.
		
	\end{enumerate}
	
	% ============================================================
	\section{Kế hoạch thực nghiệm}
	
	Nghiên cứu được triển khai theo các giai đoạn sau.
	
	\begin{table}[H]
		\centering
		\caption{Lộ trình thực nghiệm.}
		\label{tab:roadmap}
		
		\begin{tabular}{cll}
			\toprule
			Giai đoạn & Nội dung & Kết quả \\
			\midrule
			1 & Reproduce AgentMove & Baseline \\
			2 & GPS preprocessing & GPS benchmark \\
			3 & Traditional baselines & Baseline results \\
			4 & LLM baseline & LLM results \\
			5 & AgentMove-GPS & Agent baseline \\
			6 & Standalone BBN & BBN baseline \\
			7 & LLM + BBN & Initial proposed model \\
			8 & Belief Memory & Sequential belief model \\
			9 & Adaptive reasoning & Final model \\
			10 & Robustness & Noise and missing GPS \\
			11 & Calibration & Uncertainty evaluation \\
			12 & Cross-city & Generalization \\
			13 & Cost analysis & Efficiency evaluation \\
			14 & Case studies & Qualitative analysis \\
			\bottomrule
		\end{tabular}
		
	\end{table}
	
	% ============================================================
	\section{Thiết kế thí nghiệm tổng thể}
	
	Một benchmark hoàn chỉnh có thể được tổ chức như sau:
	
	\begin{equation}
		Datasets
		\times
		Models
		\times
		Noise
		\times
		MissingRate
		\times
		Seeds.
	\end{equation}
	
	Trong đó:
	
	\begin{equation}
		Datasets
		=
		\{
		GeoLife,
		T\text{-}Drive,
		Porto
		\},
	\end{equation}
	
	\begin{equation}
		Models
		=
		\{
		LSTM,
		Transformer,
		LLM,
		AgentMove,
		AgentMove\text{-}BBN,
		BeliefMove
		\}.
	\end{equation}
	
	% ============================================================
	\section{Phân tích trường hợp điển hình}
	
	Một trajectory có thể có dạng:
	
	\begin{equation}
		P_1
		\rightarrow
		P_2
		\rightarrow
		P_3
		\rightarrow
		P_4.
	\end{equation}
	
	Tại $P_4$, Agent sinh:
	
	\begin{equation}
		P_{LLM}
		=
		\{
		Road_A:0.72,
		Road_B:0.18,
		Road_C:0.10
		\}.
	\end{equation}
	
	Trong khi BBN có:
	
	\begin{equation}
		P_{BBN}
		=
		\{
		Road_A:0.41,
		Road_B:0.49,
		Road_C:0.10
		\}.
	\end{equation}
	
	Mức độ bất đồng được tính:
	
	\begin{equation}
		D_{KL}
		=
		D_{KL}
		(
		P_{LLM}
		\|
		P_{BBN}
		).
	\end{equation}
	
	Nếu mức độ bất đồng cao, Agent kích hoạt thêm reasoning.
	
	Sau khi có thêm context:
	
	\begin{equation}
		P^*
		=
		\{
		Road_A:0.56,
		Road_B:0.34,
		Road_C:0.10
		\}.
	\end{equation}
	
	Điểm GPS tiếp theo được sinh từ belief cuối cùng.
	
	% ============================================================
	\section{Định hướng mở rộng}
	
	Framework có thể được mở rộng theo bốn hướng.
	
	\subsection{Multi-step forecasting}
	
	\begin{equation}
		p_{t+1:t+K}.
	\end{equation}
	
	\subsection{Destination prediction}
	
	Thay vì chỉ dự đoán điểm kế tiếp:
	
	\begin{equation}
		P(Destination\mid\mathcal{T}).
	\end{equation}
	
	\subsection{Trajectory generation}
	
	Sinh toàn bộ trajectory có xác suất:
	
	\begin{equation}
		P(
		p_{t+1},\ldots,p_{t+K}
		\mid
		\mathcal{T}_{1:t}
		).
	\end{equation}
	
	\subsection{Real-time mobility agent}
	
	Agent có thể liên tục nhận GPS stream:
	
	\begin{equation}
		GPS_t
		\rightarrow
		Belief_t
		\rightarrow
		Prediction_t
		\rightarrow
		GPS_{t+1}
		\rightarrow
		Belief_{t+1}.
	\end{equation}
	
	Điều này mở ra khả năng ứng dụng trong Intelligent Transportation Systems.
	
	% ============================================================
	\section{Kết luận}
	
	Nghiên cứu đề xuất BeliefMove, một framework kết hợp Large Language
	Model Agent, Bayesian Belief Network và adaptive reasoning cho bài toán
	dự đoán điểm GPS kế tiếp. Ý tưởng trung tâm là chuyển bài toán từ
	deterministic next-point prediction sang probabilistic agentic mobility
	reasoning.
	
	Trong framework đề xuất, LLM không trực tiếp quyết định một điểm GPS
	duy nhất mà đóng vai trò semantic reasoner, phân tích trajectory, suy
	luận mobility intent và sinh các giả thuyết di chuyển. Bayesian Belief
	Network đóng vai trò probabilistic reasoner, định lượng xác suất của các
	giả thuyết và cập nhật belief khi quan sát mới xuất hiện. Agent đóng vai
	trò decision maker, sử dụng uncertainty và disagreement để quyết định
	khi nào cần thực hiện thêm reasoning.
	
	Mô hình vì vậy được xây dựng dựa trên ba nguyên lý:
	
	\begin{equation}
		\boxed{
			Semantic\ Reasoning
			+
			Probabilistic\ Belief
			+
			Adaptive\ Decision
		}
	\end{equation}
	
	Khung nghiên cứu này không chỉ hướng tới cải thiện ADE hoặc FDE mà còn
	đặt trọng tâm vào robustness, uncertainty calibration, cross-city
	generalization và inference efficiency. Đây là các yếu tố cần thiết để
	LLM-based trajectory prediction có thể tiến từ một mô hình dự đoán vị
	trí sang một hệ thống agentic mobility intelligence có khả năng suy luận
	và ra quyết định dưới điều kiện bất định.
	
\end{document}